% CodeAuthorityBench: A Benchmark for Authority-Aware Code Retrieval in AI Coding Agents
% Target: arXiv preprint / MSR Data & Tool track (ACM acmart, sigconf)

\documentclass[sigconf,nonacm]{acmart}

\usepackage{booktabs}
\usepackage{multirow}
\usepackage{graphicx}
\usepackage{amsmath}
\usepackage{xcolor}

% Remove ACM reference/copyright for preprint
\setcopyright{none}
\settopmatter{printacmref=false}
\renewcommand\footnotetextcopyrightpermission[1]{}
\pagestyle{plain}

\title{CodeAuthorityBench: A Benchmark for Authority-Aware\\Code Retrieval in AI Coding Agents}

\author{Sahil Soni}
\affiliation{%
  \institution{Independent Researcher}
}
\email{sahil.soni2409@gmail.com}

\begin{abstract}
AI coding agents rely on retrieved repository context to explain, modify, and debug software systems.
Existing retrieval methods optimize for textual or semantic relevance, but relevance alone does not identify which code governs system behavior.
In large repositories, database schemas, configuration files, core domain modules, and shared contracts may be structurally more authoritative than the files most similar to a query.
We introduce \textbf{CodeAuthorityBench}, a benchmark of 388 repository-understanding tasks across 7 real-world open-source projects, where the correct answer requires identifying source-of-truth files rather than semantically similar snippets.
We evaluate \textbf{CodeBrain}, an authority-aware re-ranking method that combines hybrid retrieval with graph-based structural authority scores derived from tree-sitter AST parsing, dependency analysis, and repository-level signals.
Our results show that authority-aware re-ranking significantly outperforms lexical, symbolic, graph-only, repo-map, and BM25-rerank baselines (all $p < 0.0001$, Wilcoxon signed-rank, large effect sizes), while remaining competitive with strong neural and hybrid retrieval.
CodeBrain leads 7 of 10 metrics, including all top-$k$ recall, NDCG, authority recall, and blast-radius F1.
We release CodeAuthorityBench, the CodeBrain tool as an open-source MCP server, and full reproduction scripts.
\end{abstract}

\keywords{code retrieval, authority scoring, AI coding agents, benchmark, structural analysis}

\maketitle

%% ============================================================
\section{Introduction}
%% ============================================================

AI coding agents---tools such as Claude Code, Cursor, and GitHub Copilot---increasingly rely on retrieved repository context to explain code, plan modifications, and predict the impact of changes.
The quality of retrieved context directly determines whether an agent's edits are safe or destructive.

Existing retrieval methods for code span lexical search (BM25), dense embedding retrieval, hybrid combinations, symbol-level matching, and call-graph traversal.
These methods optimize for \emph{relevance}: returning files whose content or symbols best match a query.
However, relevance is not authority.
In a large codebase, a database schema file may be the source of truth for data validation, yet rank below a verbose test fixture that mentions the same field names.
A shared configuration module may govern the behavior of 40 downstream services, yet appear less ``relevant'' than a single endpoint handler that references the same terms.

We define \textbf{code authority} as a structural property: a file's authority is determined by its position in the dependency graph, its role as a schema or interface definer, its import centrality, and its stability over time.
Authority-aware retrieval aims to surface files that \emph{govern} the queried behavior, not merely files that \emph{mention} it.

This work makes four contributions:

\begin{enumerate}
    \item \textbf{Problem definition.} We formalize authority-aware code retrieval as a distinct retrieval objective from relevance-based code search.
    \item \textbf{Benchmark.} We release CodeAuthorityBench, a dataset of 388 tasks across 7 real-world repositories (Express.js, FastAPI, Django, Flask, Gin, NestJS, Spring Framework), with four task types: source-of-truth identification, schema location, blast-radius prediction, and edit-risk classification.
    \item \textbf{Method.} We introduce CodeBrain, an authority-aware re-ranking layer that combines hybrid (BM25 + dense) retrieval with deterministic structural authority scores computed from tree-sitter AST parsing, dependency graphs, and repository-level signals.
    \item \textbf{Empirical findings.} Authority-aware re-ranking significantly improves retrieval over all non-neural baselines ($p < 0.0001$, large effect sizes) and adds consistent gains on structural metrics (authority recall, blast-radius F1, NDCG) over strong neural retrieval, while remaining competitive on MRR and R@1.
\end{enumerate}

Our work is motivated by TreeBench~\cite{treebench2026}, which showed empirical evidence that flat retrieval struggles when hierarchy determines answer authority in regulatory document structures.
CodeAuthorityBench tests whether the same structural failure exists in software repositories.


%% ============================================================
\section{Related Work}
%% ============================================================

\paragraph{Code retrieval for AI agents.}
Code search has evolved from keyword matching to semantic retrieval using embedding models~\cite{feng2020codebert,guo2022unixcoder}.
Tools such as GitHub code search, Sourcegraph, and Aider's repo-map~\cite{aider2024} provide various retrieval strategies.
Many tools still rely heavily on flat semantic or symbol retrieval and do not explicitly model structural authority.

\paragraph{Structural code analysis.}
AST-based and call-graph-based tools have a long history in software engineering~\cite{reps1998program}.
Tree-sitter provides fast incremental parsing across languages.
Tools like Understand, Sourcetrail, and ctags extract structural information, but their output is typically used for navigation rather than retrieval ranking.

\paragraph{Hierarchical retrieval.}
TreeBench~\cite{treebench2026} demonstrated that flat retrieval achieves only 45\% required-node recall on hierarchical regulatory structures, even when answer accuracy reaches 76\%.
This gap motivated the hypothesis that similar structural failures occur in code retrieval, where import hierarchies, dependency graphs, and schema ownership create implicit authority relationships.

\paragraph{Repo-level code understanding.}
Recent work on repository-level code understanding~\cite{zhang2023repocoder,liu2023repobench} focuses on cross-file context for code completion.
These approaches optimize for predicting the next code token rather than identifying which files govern a behavior---a distinct retrieval objective that CodeAuthorityBench evaluates.


%% ============================================================
\section{Code Authority}
%% ============================================================

We define \textbf{code authority} as a graph-based property of files within a repository.
A file's authority score reflects its structural control over downstream behavior, independent of query similarity.

\subsection{Authority Scoring}

Given a repository $R$ with files $F = \{f_1, \ldots, f_n\}$ and a dependency graph $G = (F, E)$ where $(f_i, f_j) \in E$ indicates $f_i$ imports or calls $f_j$, the authority score is:

\begin{equation}
\text{Auth}(f) = \sum_{s \in S} w_s \cdot \sigma_s(f, G, R)
\end{equation}

where $S$ is the set of structural signals and $\sigma_s$ is the normalized signal value:

\begin{itemize}
    \item \textbf{Dependency centrality} ($w_1 = 0.35$): Normalized in-degree in the import/call graph. Files that many modules depend on score higher.
    \item \textbf{Schema/interface ownership} ($w_2 = 0.20$): Binary detection of type definitions, interface declarations, schema files, and configuration exports.
    \item \textbf{Churn stability} ($w_3 = 0.15$): Inverse of normalized commit frequency. Stable files that rarely change score higher.
    \item \textbf{Directory prior} ($w_4 = 0.15$): Heuristic based on directory path. Files in \texttt{core/}, \texttt{schema/}, \texttt{config/} directories receive a bonus; files in \texttt{test/}, \texttt{scripts/}, \texttt{examples/} receive a penalty.
    \item \textbf{Out-degree penalty} ($w_5 = 0.15$): Normalized out-degree penalty. Files that import many other modules are consumers, not providers.
\end{itemize}

All signals are deterministic: the same repository and configuration always produce the same scores.
No LLM, no embeddings, no randomness between runs.

\subsection{Authority-Aware Re-ranking}

CodeBrain operates as a re-ranking layer on top of existing retrieval:

\begin{enumerate}
    \item A hybrid retrieval system (BM25 + dense embeddings) generates a candidate pool of 50 files per query.
    \item Each candidate receives a hybrid relevance score: $\text{rel}(f) = 0.5 \cdot \text{BM25}(f) + 0.5 \cdot \text{Dense}(f)$.
    \item The final score combines relevance with authority: $\text{score}(f) = (1 - \alpha) \cdot \text{rel}(f) + \alpha \cdot \text{Auth}(f)$, where $\alpha = 0.3$.
    \item Files are re-ranked by the combined score.
\end{enumerate}

This architecture means CodeBrain adds an authority layer on top of retrieval rather than replacing it.
The authority weight $\alpha = 0.3$ was selected to preserve retrieval quality while boosting structurally important files.


%% ============================================================
\section{CodeAuthorityBench}
%% ============================================================

CodeAuthorityBench is a benchmark of repository-understanding tasks where the correct answer requires identifying source-of-truth files, not just semantically similar snippets.

\subsection{Repositories}

We selected 7 widely-used open-source repositories spanning 5 programming languages (JavaScript, TypeScript, Python, Go, and Java) and diverse architectural patterns (Table~\ref{tab:repos}).

\begin{table}[h]
\centering
\caption{Target repositories in CodeAuthorityBench.}
\label{tab:repos}
\begin{tabular}{llr}
\toprule
Repository & Language & Files \\
\midrule
Express.js & JavaScript & 141 \\
FastAPI & Python & 1,133 \\
Django & Python & 2,967 \\
Flask & Python & 83 \\
Gin & Go & 99 \\
NestJS & TypeScript & 1,683 \\
Spring Framework & Java & 4,761 \\
\bottomrule
\end{tabular}
\end{table}

\subsection{Task Types}

Each task consists of a natural-language query, a task type, and a set of ground-truth authoritative files annotated by developer judgment.
We define four task types:

\begin{itemize}
    \item \textbf{what\_governs} (29\%): ``What file controls X?'' Requires identifying the source-of-truth file for a behavior.
    \item \textbf{where\_is\_schema} (25\%): ``Where is the schema for X?'' Requires locating type definitions, interfaces, or data models.
    \item \textbf{what\_breaks} (25\%): ``What breaks if X changes?'' Requires predicting downstream impact (blast radius).
    \item \textbf{is\_safe\_to\_modify} (21\%): ``Is it safe to modify X?'' Requires classifying edit risk as SAFE, CAUTION, or REQUIRES\_HUMAN\_APPROVAL.
\end{itemize}

The benchmark contains 388 tasks total: approximately 55 per repository, with a controlled 8-task fixture for sanity checking.
Ground-truth files are annotated based on structural authority---dependency centrality, schema ownership, and behavioral control---not textual similarity to the query.

\subsection{Annotation}

Tasks were generated using LLM-assisted annotation and validated against repository structure.
Each task includes a rationale explaining why the ground-truth files are authoritative.
Annotation guidelines specify that ground-truth files should be those with structural authority over the queried behavior, prioritizing files that other modules depend on over files that merely mention the queried concept.


%% ============================================================
\section{Experiments}
%% ============================================================

\subsection{Baselines}

We evaluate 7 baselines spanning lexical, symbolic, graph-based, and neural retrieval:

\begin{itemize}
    \item \textbf{BM25}: Okapi BM25 over raw file content.
    \item \textbf{SymbolSearch}: Regex-based extraction of function/class/variable names, BM25 over extracted symbols.
    \item \textbf{CallGraph}: Regex-based call-graph construction, ranking by call-graph distance from query-matched files.
    \item \textbf{RepoMap}: Aider-style structured summaries (file path + class/function names), BM25 over summaries.
    \item \textbf{BM25+Rerank}: BM25 candidates re-ranked by authority heuristics (no graph).
    \item \textbf{DenseRetrieval}: Dense embedding retrieval using Gemini text-embedding-004.
    \item \textbf{Hybrid}: BM25 + DenseRetrieval with score fusion ($\beta = 0.5$).
\end{itemize}

\textbf{CodeBrain} uses the same Hybrid candidate pool, re-ranked with real tree-sitter authority scores ($\alpha = 0.3$).

\subsection{Metrics}

We report 10 metrics covering retrieval quality, authority identification, and downstream impact prediction:

\begin{itemize}
    \item \textbf{Retrieval}: P@5, R@1, R@5, R@10, MRR, NDCG@5, NDCG@10
    \item \textbf{Authority}: Authority Recall (binary: did any ground-truth file appear in top-$k$?)
    \item \textbf{Edit risk}: An exploratory retrieval-based proxy for \texttt{is\_safe\_to\_modify} tasks, scored as whether the top-ranked file is among the ground-truth authoritative files (binary top-1). This is a coarse proxy rather than a full risk-tier (SAFE/CAUTION/REQUIRES\_HUMAN\_APPROVAL) classification, which we leave to future work
    \item \textbf{Blast radius}: F1 between predicted and actual affected files, evaluated at $k = |\text{ground truth}|$ to avoid fixed-$k$ precision penalties
\end{itemize}

Statistical significance is assessed via Wilcoxon signed-rank tests on per-task MRR with continuity correction.
95\% confidence intervals are computed via bootstrap (1,000 samples, seeded RNG).

\subsection{Results}

Table~\ref{tab:main} presents the overall results across all 388 tasks and 8 methods.

\begin{table*}[t]
\centering
\caption{Overall results on CodeAuthorityBench (388 tasks, 7 repositories). Best in \textbf{bold}. CodeBrain leads 7 of 10 metrics.}
\label{tab:main}
\begin{tabular}{lcccccccccc}
\toprule
Method & P@5 & R@1 & R@5 & R@10 & MRR & N@5 & N@10 & Auth & EditR & BRF1 \\
\midrule
BM25           & 12.7 & 16.1 & 42.5 & 53.5 & 35.9 & 33.1 & 37.4 & 65.5 & 19.0 & 17.3 \\
SymbolSearch   &  8.8 &  6.9 & 25.4 & 37.8 & 21.7 & 18.8 & 23.4 & 48.2 & 19.7 & 15.8 \\
CallGraph      &  1.5 &  0.8 &  3.9 &  8.3 &  3.8 &  2.9 &  4.5 & 11.6 & 17.6 &  2.7 \\
RepoMap        & 13.9 & 20.7 & 41.3 & 48.0 & 39.8 & 34.5 & 37.0 & 62.1 & 21.1 & 22.6 \\
BM25+Rerank    & 13.9 & 17.2 & 44.2 & 56.8 & 38.1 & 34.8 & 39.5 & 68.8 & 20.4 & 19.0 \\
DenseRetrieval & 16.9 & \textbf{30.7} & 56.0 & 66.7 & \textbf{55.1} & 48.5 & 52.8 & 81.2 & \textbf{23.9} & 24.9 \\
Hybrid         & 18.0 & 28.7 & 58.0 & 68.6 & 54.0 & 48.9 & 53.1 & 82.2 & 22.5 & 27.9 \\
\textbf{CodeBrain} & \textbf{18.9} & 27.3 & \textbf{60.5} & \textbf{71.3} & 54.3 & \textbf{49.8} & \textbf{54.1} & \textbf{83.8} & 22.5 & \textbf{28.4} \\
\bottomrule
\end{tabular}
\end{table*}

CodeBrain significantly outperforms lexical, symbolic, graph-only, repo-map, and BM25-rerank baselines, while remaining competitive with strong neural and hybrid retrieval.
Its clearest gains are in authority recall, top-$k$ recall, NDCG, and blast-radius estimation.

\subsection{Statistical Significance}

Table~\ref{tab:significance} presents Wilcoxon signed-rank tests comparing CodeBrain against each baseline on per-task MRR.

\begin{table}[h]
\centering
\caption{Wilcoxon signed-rank tests (CodeBrain vs.\ baselines, per-task MRR). $\Delta$MRR is the mean per-task MRR difference over \emph{non-tied} pairs; because the Wilcoxon test excludes zero-difference (tied) pairs, the number of contributing pairs is $n < 388$ and $\Delta$MRR therefore differs from the difference of aggregate MRR values in Table~\ref{tab:main}. $^{***}p < 0.001$.}
\label{tab:significance}
\begin{tabular}{lrrrr}
\toprule
Baseline & $Z$ & $p$ & $r$ & $\Delta$MRR \\
\midrule
BM25           &  8.78 & $<$0.0001$^{***}$ & 0.55 & +0.281 \\
SymbolSearch   & 11.53 & $<$0.0001$^{***}$ & 0.66 & +0.419 \\
CallGraph      & 14.93 & $<$0.0001$^{***}$ & 0.83 & +0.613 \\
RepoMap        &  5.16 & $<$0.0001$^{***}$ & 0.31 & +0.198 \\
BM25+Rerank    &  8.24 & $<$0.0001$^{***}$ & 0.54 & +0.266 \\
DenseRetrieval &  0.31 & 0.694 & 0.02 & $-$0.011 \\
Hybrid         &  0.16 & 0.827 & 0.01 & +0.008 \\
\bottomrule
\end{tabular}
\end{table}

The results reveal two distinct regimes.
Against non-neural baselines (BM25, SymbolSearch, CallGraph, RepoMap, BM25+Rerank), CodeBrain achieves highly significant improvements with medium-to-large effect sizes ($r = 0.31$--$0.83$).
Against neural baselines (DenseRetrieval, Hybrid), the differences on MRR are not statistically significant.
However, CodeBrain consistently leads on structural metrics: authority recall (83.8\% vs 82.2\% Hybrid, 81.2\% Dense), NDCG@10 (54.1\% vs 53.1\%, 52.8\%), and blast-radius F1 (28.4\% vs 27.9\%, 24.9\%).

This pattern suggests that authority re-ranking is \emph{complementary} to neural retrieval: it adds the most value when the base retrieval lacks structural awareness, and provides consistent (though smaller) gains even when combined with strong neural retrieval.

\subsection{Per-Task-Type Analysis}

Table~\ref{tab:pertask} shows CodeBrain's advantage varies by task type.

\begin{table}[h]
\centering
\caption{CodeBrain vs.\ Hybrid by task type (P@5 / Auth Recall).}
\label{tab:pertask}
\begin{tabular}{lcccc}
\toprule
Task Type & Hybrid P@5 & CB P@5 & Hybrid Auth & CB Auth \\
\midrule
what\_governs    & 6.1 & \textbf{6.9} & 28.9 & \textbf{29.1} \\
where\_is\_schema & 4.6 & \textbf{4.6} & 22.9 & \textbf{22.9} \\
what\_breaks     & 4.3 & \textbf{4.3} & 16.0 & \textbf{17.3} \\
is\_safe\_to\_modify & 3.0 & \textbf{3.0} & 14.4 & \textbf{14.4} \\
\bottomrule
\end{tabular}
\end{table}

The largest gains appear on \texttt{what\_governs} tasks---precisely the task type where identifying the structural source of truth matters most---and on \texttt{what\_breaks} tasks where blast-radius prediction benefits from authority-aware ranking.

\subsection{Ablation Study}

To measure the contribution of individual authority signals, we run ablation variants that selectively enable signal subsets (Table~\ref{tab:ablation}).
All variants use the same hybrid retrieval; only the authority re-ranking signals differ.

\begin{table}[h]
\centering
\caption{Ablation study: contribution of individual authority signals. Variants come from a separate configuration sweep over the same pipeline; absolute values differ slightly from Table~\ref{tab:main} (e.g., \textbf{flat-only} corresponds to the Hybrid baseline), so this table should be read for the \emph{relative} contribution of each signal rather than for absolute parity with the main results.}
\label{tab:ablation}
\begin{tabular}{lccccc}
\toprule
Variant & P@5 & R@5 & MRR & Auth \\
\midrule
flat-only (no authority) & 18.0 & 58.0 & 54.3 & 82.2 \\
directory-only           & 18.6 & 59.8 & 56.0 & 84.0 \\
graph-only               & 18.8 & 59.8 & 56.7 & 85.3 \\
graph+directory          & 19.2 & 61.0 & 57.7 & 85.6 \\
graph+schema             & 19.4 & 61.1 & 57.6 & 85.8 \\
\textbf{full}            & \textbf{19.5} & \textbf{61.7} & 56.5 & 85.6 \\
\bottomrule
\end{tabular}
\end{table}

Key observations:
\begin{itemize}
    \item \textbf{Graph signals dominate.} Graph-only outperforms directory-only on MRR (56.7 vs 56.0) and Auth Recall (85.3 vs 84.0), confirming that structural analysis adds value beyond cheap path heuristics.
    \item \textbf{Each signal adds value.} Full (all signals) achieves the best P@5 and R@5. No single signal dominates all metrics.
    \item \textbf{Graph+schema is the best single combination} for Auth Recall (85.8\%) and R@1, validating that type-level analysis adds signal beyond graph structure alone.
    \item \textbf{All authority variants beat flat retrieval} across every metric, confirming that authority re-ranking is consistently beneficial.
\end{itemize}


%% ============================================================
\section{Discussion}
%% ============================================================

\paragraph{Authority as a complementary signal.}
Our results show that authority-aware re-ranking is not a replacement for strong neural retrieval but a complementary layer.
The largest gains come from adding authority to weaker retrieval methods (BM25, SymbolSearch), while gains over Hybrid/Dense are consistent but modest.
This is expected: neural embeddings already capture some structural patterns implicitly through training data, while authority scoring makes these patterns explicit and deterministic.

\paragraph{Why neural retrieval wins on R@1 and MRR.}
DenseRetrieval achieves the highest R@1 (30.7\%) and MRR (55.1\%).
These metrics reward placing a single correct file at position 1.
Dense embeddings excel at semantic matching, which often identifies the most obviously relevant file.
Authority re-ranking may sometimes promote a \emph{more authoritative} file above the \emph{most semantically similar} one, which improves precision and recall at deeper cutoffs but can lower R@1 when the promoted file is authoritative but not in the ground truth.

\paragraph{Practical implications.}
For AI coding agents, the most dangerous errors come not from missing a relevant file but from modifying a structurally critical file without understanding its authority.
Authority recall (83.8\%) and blast-radius F1 (28.4\%) directly measure this safety-relevant capability.
CodeBrain's consistent advantage on these metrics suggests practical value for agent governance, even when overall MRR differences are small.

\paragraph{Error analysis: foundation vs.\ orchestration authority.}
To probe what the authority score captures beyond the benchmark, we applied CodeBrain to three non-benchmark repositories spanning two languages: an independent 189-file Python operations agent, CodeBrain's own 28-module TypeScript implementation, and LangGraph, a 453-file Python agent-orchestration monorepo.
Across all three, the model reliably elevated \emph{foundation} files---those that define shared state, schemas, persistence interfaces, and common error and type definitions---to the ROOT tier (in the Python agent, the connection layer, repository modules, and shared exceptions scored highest).
This supports interpreting authority as structural \emph{blast radius} rather than human-perceived importance.
The contrast is sharpest in LangGraph: \texttt{graph/state.py}, which defines the central \texttt{StateGraph} abstraction, ranked first, while \texttt{pregel/main.py}, the Pregel execution engine, ranked 67th of 453.
More broadly, entry-point and orchestration files were consistently ranked \emph{below} expectation: in the Python agent, the command-line entry point and application \texttt{\_\_main\_\_} module fell into the LEAF tier (ranked 85th and 48th of 189 files), and the TypeScript implementation's server entry points ranked last.
These files coordinate runtime execution but are dependency \emph{sinks}---they import many modules and are imported by none---so an in-degree-driven score treats them as disposable.
This exposes a genuine distinction: repository authority is multi-dimensional.
\emph{Foundation authority}, which the current model measures through dependency and contract centrality, captures files that many components depend on; this is the correct signal for the \texttt{guard\_change} and \texttt{is\_safe\_to\_modify} use cases, since a dependency sink genuinely has near-zero blast radius.
\emph{Orchestration authority}---files that initiate or coordinate behavior---is orthogonal and is not captured by the current signals, so orchestrators may be important to developers yet under-ranked for behavior-governance queries.
Conversely, widely imported utility files, such as diagramming helpers, can be over-ranked.
We therefore treat CodeBrain v1 as a foundation-authority baseline and leave explicit orchestration/entry-point authority modeling to future work, evaluated as an optional ablation rather than folded into the reported configuration to avoid post-hoc tuning of the benchmark results.

\paragraph{Limitations.}
(1) Ground-truth annotations are LLM-assisted and may contain noise; inter-annotator agreement was not measured.
(2) The authority weight $\alpha = 0.3$ was not extensively tuned; systematic hyperparameter search may yield better configurations.
(3) CodeAuthorityBench covers 7 repositories in 5 languages; generalization to other languages and repository architectures requires further evaluation.
(4) The benchmark evaluates static retrieval; integration with actual agent workflows (e.g., measuring whether authority-aware retrieval reduces agent edit errors) is left to future work.
(5) As the error analysis above indicates, the current model captures foundation authority but not orchestration authority; entry-point and coordinator files may therefore be under-ranked for behavior-governance queries even when correctly ranked for blast-radius and safe-to-modify queries.
(6) \textbf{Label--signal overlap (construct validity).} Ground-truth authority labels were annotated using structural criteria---dependency centrality, schema ownership, and behavioral control---that overlap with the signals CodeBrain scores. The benchmark therefore currently measures how faithfully the method operationalizes this structural definition of authority, rather than validating that definition against an independent notion of authority. Re-annotation using criteria independent of the scoring signals, together with inter-annotator agreement (Cohen's~$\kappa$) from external annotators, is pending; we release an annotation kit to support this. Results should be read with this overlap in mind.
(7) Several benchmark repositories (e.g., Flask, Django, Express) are widely represented in the pretraining data of the neural baselines, which may inflate their absolute scores; evaluating on lower-profile repositories is future work.
(8) At the evaluated commits, five ground-truth paths (three in Express, two in FastAPI) do not resolve to files present in the checkout and were scored as misses; these are documented in the released \texttt{PIN\_PROVENANCE.md} and affect five tasks, not the aggregate conclusions.


%% ============================================================
\section{Conclusion}
%% ============================================================

We introduced CodeAuthorityBench, a benchmark of 388 repository-understanding tasks designed to evaluate whether retrieval methods can identify structurally authoritative code, not just semantically similar files.
We presented CodeBrain, an authority-aware re-ranking method that combines hybrid retrieval with deterministic structural authority scores.

Our experiments show that authority-aware re-ranking significantly outperforms lexical, symbolic, graph-only, repo-map, and BM25-rerank baselines, while remaining competitive with strong neural and hybrid retrieval.
CodeBrain leads 7 of 10 metrics, with its clearest advantages in authority recall, NDCG, and blast-radius estimation---metrics most relevant to the safety of AI coding agent behavior.

We release CodeAuthorityBench, the CodeBrain MCP server, and full reproduction scripts to support further research on authority-aware code retrieval.


%% ============================================================
\section*{Data and Code Availability}
%% ============================================================

CodeAuthorityBench---all 388 tasks with ground-truth annotations and rationales---together with the CodeBrain implementation and MCP server, the seven baseline implementations, and the complete evaluation and statistical-analysis scripts, is publicly available at \url{https://github.com/whitepaper27/codebrain}.
For reproducibility, each benchmark repository is pinned to the exact commit it was evaluated against (documented in \texttt{benchmarks/PIN\_PROVENANCE.md}), and the released scripts regenerate every table in this paper from those pinned states.
The software and dataset are released under the Apache~2.0 license.


%% ============================================================
%% References
%% ============================================================

\bibliographystyle{ACM-Reference-Format}

\begin{thebibliography}{10}

\bibitem{treebench2026}
Sahil Soni. 2026.
\newblock TreeBench: Evaluating Hierarchical Retrieval for Regulatory Question Answering.
\newblock Manuscript. \url{https://github.com/whitepaper27/TreeBench}.

\bibitem{feng2020codebert}
Zhangyin Feng, Daya Guo, Duyu Tang, Nan Duan, Xiaocheng Feng, Ming Gong, Linjun Shou, Bing Qin, Ting Liu, Daxin Jiang, and Ming Zhou. 2020.
\newblock CodeBERT: A Pre-Trained Model for Programming and Natural Languages.
\newblock \emph{Findings of EMNLP 2020}.

\bibitem{guo2022unixcoder}
Daya Guo, Shuai Lu, Nan Duan, Yanlin Wang, Ming Zhou, and Jian Yin. 2022.
\newblock UniXcoder: Unified Cross-Modal Pre-training for Code Representation.
\newblock \emph{ACL 2022}.

\bibitem{aider2024}
Paul Gauthier. 2024.
\newblock Aider: AI pair programming in your terminal.
\newblock \url{https://aider.chat}.

\bibitem{reps1998program}
Thomas Reps. 1998.
\newblock Program analysis via graph reachability.
\newblock \emph{Information and Software Technology} 40, 11-12 (1998), 701--726.

\bibitem{zhang2023repocoder}
Fengji Zhang, Bei Chen, Yue Zhang, Jin Liu, Daoguang Zan, Yi Mao, Jian-Guang Lou, and Weizhu Chen. 2023.
\newblock RepoCoder: Repository-Level Code Completion Through Iterative Retrieval and Generation.
\newblock \emph{EMNLP 2023}.

\bibitem{liu2023repobench}
Tianyang Liu, Canwen Xu, and Julian McAuley. 2023.
\newblock RepoBench: Benchmarking Repository-Level Code Auto-Completion Systems.
\newblock \emph{arXiv preprint arXiv:2306.03091}.

\end{thebibliography}

\end{document}
